% ECONOMICS_PROBLEM: {"id": "L42-541", "title": "Vertical integration and self-preferencing", "jel_code": "L42", "source_id": "OP-0077"}
% !TeX program = xelatex
\documentclass[11pt,a4paper]{article}
\usepackage[margin=23mm]{geometry}
\usepackage{fontspec}
\setmainfont{Latin Modern Roman}
\usepackage{amsmath,amssymb,mathtools}
\usepackage{xcolor,enumitem,booktabs,longtable,array,xurl,hyperref}
\definecolor{muted}{HTML}{53636C}
\hypersetup{colorlinks=true,urlcolor=blue,linkcolor=blue}
\setlength{\parindent}{0pt}
\setlength{\parskip}{5pt}
\newcommand{\R}{\mathbb R}
\newcommand{\E}{\mathbb E}
\newcommand{\N}{\mathbb N}
\newcommand{\ind}{\mathbf 1}
\newcommand{\Var}{\operatorname{Var}}
\newcommand{\Cov}{\operatorname{Cov}}
\newcommand{\supp}{\operatorname{supp}}
\DeclareMathOperator*{\argmin}{arg\,min}
\DeclareMathOperator*{\argmax}{arg\,max}
\newcommand{\entrymeta}[3]{{\sffamily\small\color{muted}\textbf{\texttt{#1}}\quad|\quad #2\par #3\par}}
\newcommand{\Problemref}[1]{\texttt{#1}}
\newcommand{\JELref}[2]{\texttt{#2} (JEL \texttt{#1})}
\begin{document}
\section*{Vertical integration and self-preferencing}
\textbf{Problem ID:} \texttt{L42-541}\quad \textbf{JEL:} \texttt{L42}\par
% BEGIN ECONOMICS STATEMENT
\label{problem:OP-0077}
\label{atlas:prob:I7}

\noindent\textbf{Original identifier:} I7 (atlas item 77). \textbf{Classification:} Vertical conduct and welfare research program. \textbf{Original tractability label:} Medium (editorial, not a complexity result).

\subsection*{Definitions and assumptions}

Let an upstream firm supply an input to downstream firms, one of which may be integrated with it. A regime $v$ specifies ownership, wholesale contracting, access obligations, ranking rules, and permissible discrimination. Demand depends on consumer prices, quality, and accessibility. Supply behavior follows a declared bargaining or contracting game, including whether contracts are public, enforceable, and contingent. Let $e_\sigma(v)$ denote an equilibrium under a declared selection rule. Define foreclosure as a specified reduction in rivals' access, raising of their relevant costs, or exclusion relative to the baseline; do not identify it merely with a rival's lower profit.

\subsection*{Formal research question}

Identify or bound total welfare effects $\Delta W^\sigma(v,v_0)$ and consumer surplus effects, and separate mechanisms where the counterfactuals are well defined. For example, define $v_E$ as allowing internal elimination of successive price distortions while requiring the baseline access terms, and compare
\[
W(e_\sigma(v_E))-W(e_\sigma(v_0)),\qquad
W(e_\sigma(v))-W(e_\sigma(v_E)).
\]
Determine conditions under which both comparisons are identified and their signs robust across admissible conduct and bargaining assumptions. Evaluate specific equal-access or ranking restrictions using causal policy variation or a validated structural model.

\subsection*{Interpretation and scope}

The same integration can remove successive markups, coordinate investment, and change the integrated firm's incentive to disadvantage downstream rivals. The magnitude of each channel depends on contract possibilities and bargaining power; double marginalization is not guaranteed under every nonintegrated contract. A mechanism decomposition requires feasible counterfactual regimes and is generally path-dependent when channels interact. The displayed intermediate regime is therefore an explicit modeling input, not a uniquely natural attribution. Search self-preferencing is a particular conduct change whose welfare depends on match quality, ranking incentives, and outside options, and cannot be inferred solely from vertical ownership. Data on wholesale transfers may identify neither marginal costs nor bargaining threats without additional restrictions. Equal-access mandates can alter investment or participation, so the policy outcome should include these responses over a stated horizon. A successful extension would predict both rival access and final-consumer outcomes after the same intervention, with uncertainty over alternative bargaining models and selection rules.

\subsection*{Source audit and status}

[S1] develops a foreclosure mechanism; paper authors Hart and Tirole are distinguished from discussion contributors in the Brookings record. [S2] surveys foreclosure theories, and [S3] quantifies integration in television markets. All original citations are verified, but television evidence is not direct proof about platform ranking. Existing theory already characterizes efficiency and foreclosure in particular models. The explicit regime comparison is an editorial extension directed at identification and application-specific welfare assessment.

\subsection*{References}

\begin{enumerate}[label={[S\arabic*]},leftmargin=*]

\item Oliver Hart and Jean Tirole (1990). \emph{Vertical Integration and Market Foreclosure}. Brookings Papers on Economic Activity: Microeconomics, 205--286 (article and discussion). \url{https://www.brookings.edu/articles/vertical-intergration-and-market-foreclosure/}\par \textit{Source role:} Formal foreclosure mechanism; discussion authors are not paper coauthors. \textit{Locator:} Abstract and article metadata.

\item Patrick Rey and Jean Tirole (2007). \emph{A Primer on Foreclosure}. Handbook of Industrial Organization, Volume 3, Chapter 33, 2145--2220, Elsevier. \url{https://www.sciencedirect.com/science/chapter/handbook/pii/S1573448X06030330}\par \textit{Source role:} Review of foreclosure theories. \textit{Locator:} Abstract and article metadata.

\item Gregory S. Crawford, Robin S. Lee, Michael D. Whinston, and Ali Yurukoglu (2018). \emph{The Welfare Effects of Vertical Integration in Multichannel Television Markets}. Econometrica 86(3), 891--954. \url{https://doi.org/10.3982/ECTA14031}\par \textit{Source role:} Structural quantification of vertical efficiencies and foreclosure. \textit{Locator:} Abstract and article metadata.

\end{enumerate}

\subsection*{Common conventions: Source mathematical conventions}

Unless an entry states otherwise, agent and firm sets are finite. State and action spaces are measurable, strategies and outcome maps are measurable, probability laws are specified, and expectations exist for the targets under discussion. Infinite-horizon discount factors lie in $(0,1)$ and discounted objectives are finite. Norms, tolerances and complexity input sizes must be specified for computational comparisons. Constants or primitives that appear locally are inputs, not empirical facts asserted by this catalogue.

\textbf{Identification.} A statistical model $m\in\mathcal M$ implies an observable law $P_m$ and target $\theta(m)$. At law $P$, its identified set is
\[
\Theta(P;\mathcal M)=\{\theta(m):m\in\mathcal M,\ P_m=P\}.
\]
Point identification holds when this set is a singleton, or equivalently when $P_{m_1}=P_{m_2}$ implies $\theta(m_1)=\theta(m_2)$ for all compatible models. Sharp bounds mean bounds attained, or approached when the boundary is not attained, by compatible models. Writing this set is a definition, not a solution: the substantive task is to establish informative restrictions and derive or estimate the set. An unrestricted-confounding model may give only trivial bounds.

\textbf{Inference.} Identification, estimability and valid uncertainty quantification are separate questions. Fix $\alpha\in(0,1)$. A confidence procedure $C_n$ over a declared law class $\mathcal P$ has asymptotic uniform coverage at level $1-\alpha$ when
\[
\liminf_{n\to\infty}\inf_{P\in\mathcal P}\Pr_P\{\theta(P)\in C_n\}\geq1-\alpha.
\]
For set-identified targets the coverage event must be defined separately, for example coverage of the full identified set. If an entry uses identification-indexed models instead of uniquely indexed laws, the same coverage convention is applied over those models. Pointwise limits do not establish uniform validity, and asymptotic validity does not guarantee finite-sample precision. Sample sizes, dependence structures and weak-identification sequences are part of each application.

\textbf{Counterfactuals.} $Y_i(a)$ is a potential outcome under a specified intervention, including its timing, implementation and versions. It is not $Y_i$ conditional on voluntarily choosing $a$. A full intervention vector permits interference; shorthand scalar treatment notation does not silently impose no interference. Assignment, selection, attrition, anticipation and measurement must be specified. A claim about an instrument's local effect is not automatically a population policy effect.

\textbf{Economic equilibrium.} A counterfactual model includes best responses, constraints, feasibility, market clearing, state transitions and expectations consistent with the stated information. When several equilibria exist, either a declared selector is an input or the target is set-valued. Comparisons across policy regimes must use the stated selection convention. Reduced-form relationships are not presumed invariant after an equilibrium-changing intervention.

\textbf{Optimization and welfare.} Optimization is over the stated feasible set. If attainment is not established, $\sup$ or $\inf$ and $\varepsilon$-optimal choices replace an asserted maximizer or minimizer. For example, with value $V=\sup_{a\in\mathcal A}W(a)<\infty$, an $\varepsilon$-optimal set is $\{a\in\mathcal A:W(a)\geq V-\varepsilon\}$ for $\varepsilon>0$. Benefits, costs, risk and health or environmental outcomes can be aggregated only using declared units and conversion weights. Interpersonal utility weights, the treatment of fiscal transfers and foreign profits, discounting, and the behavioral welfare criterion are explicit inputs. A comparative-static derivative additionally requires existence and appropriate differentiability of the selected counterfactual.

\textbf{Sources.} Local labels [S1], [S2], and so forth restart in every question. Locators identify the relevant abstract, model, section or result at the level actually checked; a bibliographic or abstract-level verification is not represented as inspection of an entire proof. Working papers are distinguished from later publications and changing versions. Source summaries are paraphrased. All URL checks and editorial formulations refer to the audit date stated above.
% END ECONOMICS STATEMENT
\end{document}
